\documentclass[10pt,a4paper]{amsart}
\usepackage[margin=19mm]{geometry}
\usepackage{amsmath,amssymb,mathtools}
\usepackage[scr=rsfs,cal=euler]{mathalfa}
\usepackage{xfrac}
\usepackage[unicode,hidelinks,bookmarks=false]{hyperref}
\newcommand{\ie}{i.e.\ }
\newcommand{\cf}{cf.\ }
\newcommand{\resp}{respectively\ }
\newcommand{\Subsection}{\subsection}
\providecommand{\nobreakdash}{\nobreak\hbox{-}}
\setlength{\emergencystretch}{2em}
\multlinegap0pt
\usepackage{bm}
\usepackage{bbm}
\usepackage{dsfont}
\usepackage{xspace}
\usepackage{cancel}
\usepackage{mathtools}
\usepackage{amsthm}
\makeatletter
\@ifundefined{theorem}{\newtheorem{theorem}{Theorem}}{}
\@ifundefined{lemma}{\newtheorem{lemma}[theorem]{Lemma}}{}
\makeatother
\newcommand{\Par}{{\rm Par}}
\newcommand{\Stri}{{\rm S}_{_{\!\triangle\!\!}}}
\newcommand{\smtriangle}{\scalebox{.8}{$\,\triangle\,$}}
\title[Palindromicity of the numerator of a statistical generating ]{Palindromicity of the numerator of a statistical generating function}
\author{Rebecca Bourn, William Q. Erickson}
\date{Source: arXiv:2307.02652v3 (CC BY 4.0)}
\begin{document}
\maketitle

\noindent\emph{Source and licence.} Palindromicity of the numerator of a statistical generating function. Version arXiv:2307.02652v3 (CC BY 4.0), distributed under the Creative Commons Attribution 4.0 International licence, as recorded in the arXiv licence metadata for this exact version and shown on the abstract page https://arxiv.org/abs/2307.02652v3. Full rights notice: licenses/arxiv-2307.02652-CC-BY-4.0.txt.\par\smallskip

\noindent\textit{Editorial scope.} Counted unit: the Theorem whose source label is \texttt{thm:main result} in the article (source0.tex lines 349--358, source label \texttt{thm:main result}), delivered together with the article's own complete original proof of it (source0.tex lines 360--402, one \texttt{proof} environment of 43 lines). No mathematics is written, repaired or re-derived: statement and proof are the article's own text line for line. Required context retained uncounted: Wpq (lines 161--164); Npq recursion (lines 166--169); size Par ab (lines 212--215); S triangle (lines 232--235); lemma:construction (lines 247--258). Every definition, display and label of those lines is retained unchanged. Omitted from this package and not redistributed: 1--160, 165--165, 170--211, 216--231, 236--246, 259--348, 359--359, 403--637. Nothing inside a retained excerpt is omitted or altered: every retained body is delivered exactly as the article's lines are written, so no omission receipt of any kind accompanies this package. Citations: the retained excerpts carry no citation command at all, so no bibliography entry of the article is transcribed and no reference is redistributed. Reference adaptation: none was needed and none was made; every reference inside the delivered unit resolves inside it, so no unresolved reference or citation is produced. Printed slips kept literal: no printed word, symbol, hypothesis or number of the counted statement or of its proof was altered. Layout and class: the article's own class is \texttt{amsart} with packages including amssymb, ytableau, mathtools, color, amsrefs, tikz, fullpage; the wrapper is the lane's pinned \texttt{amsart} editorial wrapper at 10pt on A4, which restates the theorem-like environments the retained text needs and adds the article's own macro definitions that the retained text needs (\texttt{\textbackslash Par}, \texttt{\textbackslash Stri}, \texttt{\textbackslash smtriangle}), with extra wrapper packages (bm, bbm, dsfont, xspace, cancel, mathtools). The delivered numbering is the unit's own. Assets: no retained span contains a figure environment, a graphics-inclusion command, a PDF-page inclusion, a table or a TikZ picture; no image file is copied into this package, the package declares no asset dependency and no image file is redistributed with it. Licence: version arXiv:2307.02652v3 of this article is distributed under the Creative Commons Attribution 4.0 International licence, as recorded in the arXiv licence metadata for this exact version and shown on the abstract page https://arxiv.org/abs/2307.02652v3. A full rights notice is delivered at licenses/arxiv-2307.02652-CC-BY-4.0.txt. Characters: every retained excerpt is pure ASCII, so the delivered engine, XeLaTeX with its default Latin Modern text font, reports no missing character.
    \begin{equation}
    \label{Wpq}
    W_{pq}(t) = \sum_{k=0}^{\min\{p,q\}-1} \binom{p-1}{k} \binom{q-1}{k} t^k.
    \end{equation}

\begin{equation}
    \label{Npq recursion}
    N_{pq} = N_{p-1,q} + N_{p,q-1} - (1-t)N_{p-1,q-1} + |p-q| \cdot t \, W_{pq},
\end{equation}

\begin{equation}
    \label{size Par ab}
    |\Par(a \times b)| = \binom{a+b}{a}.
\end{equation}

\begin{equation}
\label{S triangle}
    \Stri(a,b \mid c,d) \coloneqq \sum_{\mathclap{\substack{\lambda \in \Par(a \times b), \\ \mu \in \Par(c \times d)\phantom{,}}}} |\lambda \smtriangle \mu|.
\end{equation}

\begin{lemma}
    \label{lemma:construction}
    For positive integers $k, \ell, m$, we have
    \begin{align}
    \label{lemma equation}
    \begin{split}
    \Stri(k, \ell \mid k, m) &= \Stri(k, \ell-1 \mid k, m) + \Stri(k, \ell \mid k, m-1) - \Stri(k, \ell-1 \mid k, m-1) \\
    & \hspace{10pt} + \Stri(k-1, \ell \mid k-1, m) \\
    & \hspace{+20pt} + |\ell - m| \cdot |\Par((k-1) \times \ell)|\cdot |\Par((k-1) \times m)|.
    \end{split}
    \end{align}
\end{lemma}

\begin{theorem}
    \label{thm:main result}
    Let $N_{pq}(t)$ be the polynomial defined recursively in~\eqref{Npq recursion}.
    Then
    \begin{equation}
    \label{Npq in theorem}
        N_{pq}(t) = \sum_{k=1}^{\min\{p,q\}} \Stri(k, p-k \mid k, q-k) \cdot t^k,
    \end{equation}
    where the coefficients $\Stri( - )$ are defined in~\eqref{S triangle}.
\end{theorem}

\begin{proof}
    Recall from~\eqref{Npq recursion} the recursion defining $N_{pq}(t)$, which we expand below (labeling each term for clarity later in the proof):
    \begin{equation}
        \label{Npq expanded}
        N_{pq} = \underbrace{N_{p-1,q}}_{\mathbf A} + \underbrace{N_{p,q-1}}_{\mathbf B} - \underbrace{N_{p-1, q-1}}_{\mathbf C} + \; t \cdot \underbrace{N_{p-1,q-1}}_{\mathbf D} + \; t \cdot \underbrace{|p-q| \cdot W_{pq}}_{\mathbf E},
    \end{equation}
    with initial values $N_{0,q} = N_{p,0} = N_{1,1} = 0$.
    
    We begin by verifying the initial values.
    Clearly when $p$ or $q$ is zero, the sum in~\eqref{Npq in theorem} is empty.
    If $p=q=1$, then there is a single summand where $k=1$, namely
    \[
        \Stri(1,0 \mid 1,0) \cdot t.
    \]
    Since $\Par(1 \times 0)$ contains only the empty partition, the expression above is zero, as required.

    It remains to verify the recursion itself.
    As our induction hypothesis, assume that~\eqref{Npq in theorem} is valid for all $N_{p'q'}$ with $p' < p$ or $q' < q$.
    It suffices to show, for all $k$, that the coefficient of $t^k$ in the right-hand side of~\eqref{Npq in theorem}, namely $\Stri(k, p-k \mid k, q-k)$, equals the coefficient of $t^k$ in the right-hand side of~\eqref{Npq expanded}.
    That is, writing $[t^k] f$ to denote the coefficient of $t^k$ in a polyomial $f(t)$,
    we need to show that
    \begin{align}
        \Stri(k, p-k \mid k, q-k) &= [t^k](\mathbf{A} + \mathbf{B} - \mathbf{C} + t\mathbf{D} + t\mathbf{E}) \nonumber \\
        &= [t^k]\mathbf{A} + [t^k]\mathbf{B} - [t^k]\mathbf{C} + [t^{k-1}]\mathbf{D} + [t^{k-1}]\mathbf{E}. \label{WTS} 
    \end{align}
    For ease of notation (and to align with Lemma~\ref{lemma:construction}), set $\ell \coloneqq p-k$ and $m \coloneqq q-k$.
    The required coefficients in $\mathbf{A}, \ldots, \mathbf{D}$ follow directly from~\eqref{Npq in theorem} via the induction hypothesis:
    \begin{alignat*}{3}
        [t^k]\mathbf{A} &= [t^k]N_{p-1,q} &&=\Stri(k, \ell-1 \mid k, m),\\
        [t^k]\mathbf{B} &= [t^k]N_{p,q-1} &&= \Stri(k, \ell \mid k, m-1),\\
        [t^k]\mathbf{C} &= [t^k]N_{p-1,q-1} &&= \Stri(k, \ell-1 \mid k, m-1),\\
        [t^{k-1}]\mathbf{D} &= [t^{k-1}]N_{p-1,q-1} &&= \Stri(k-1, \ell \mid k-1, m).
    \end{alignat*}
    The coefficient $[t^{k-1}]\mathbf{E}$ can be expressed by means of previous identities:
    \begin{align*}
        [t^{k-1}]\mathbf{E} &= |p-q| \cdot [t^{k-1}]W_{pq} \\
        &= |p-q| \cdot \textstyle\binom{p-1}{k-1} \binom{q-1}{k-1} && \text{by~\eqref{Wpq}}\\
        &= |(\ell+k) - (m+k)| \cdot \textstyle\binom{\ell + k-1}{k-1} \binom{m + k-1}{k-1}\\
        &= |\ell - m| \cdot |\Par((k-1) \times \ell)| \cdot |\Par((k-1) \times m)| && \text{by~\eqref{size Par ab}}.
    \end{align*}
  Upon substituting these five coefficients back into~\eqref{WTS}, we obtain the equation in Lemma~\ref{lemma:construction}, which verifies~\eqref{WTS} as desired. 
  Finally, to determine the range of the sum in~\eqref{Npq in theorem}, we observe that $\Stri(k, p-k \mid k, q-k)$ is nonzero only when $1 \leq k \leq \min\{p,q\}$.
\end{proof}

\end{document}
